
lignes n'a pu retrouver la pierre tombale de cet Albanais
d'origine. C'était un clerc assez fortuné si l'on en juge par
l'exposé de ses derniers vouloirs.

\par

On se rappelle qu'en 1484, l'église fut ruinée par les
troupes royales venues défendre la cause de Boffille de
Juge et mater la rébellion maladroite de ces pauvres
Albanais. Bien entendu, on la reconstruisit et elle dura
jusqu'à 1586 où elle fut à nouveau victime de bagarres,
cette fois-ci entre catholiques et protestants. Que reste-t-il
de cette lointaine époque ? peut-être les deux voussoirs
(ou vousseaux : pierres taillées entrant dans la construction
d'une voûte ou d'un arc), actuellement placés au-dessus
des deux piliers qui encadrent le portail d'entrée de notre
église. Ces deux pièces trouvées dans les démolitions
effectuées au cours de la transformation de la façade ouest
pendant les années 63-64 ont bien pu appartenir à un arceau
ou voûte gothique de ce temps-là, début du XVI\textsuperscript{ème}
ou fin du XV\textsuperscript{ème} siècle. Vicissitudes diverses de la malice des
temps. « Dies illi mali sunt » a dit un écrivain sacré :
« Il y eut de la grogne en ces jours-là ». (traduction
Chrysanthème).
